\documentclass{ctexbook}
\usepackage{Note}
\usepackage{unicode-math}
\setCJKmainfont[
  BoldFont = 思源宋体 Bold,
  ItalicFont = 楷体,
]{宋体}
\setcounter{secnumdepth}{3}
\begin{document}
\chapter{分子结构与光谱}
\section{分子轨道理论}
\subsection{分子的薛定谔方程}
和多电子原子类似,分子的薛定谔方程需要考虑各个核与电子本身的动能项与相互的势能项,即
\[\hat{H}\iPsi=E\iPsi,\quad\hat{H}=T_{n}+T_{e}+V_{en}+V_{ee}+V_{nn}\]
在原子单位下
\[\hat{H}=-\sum_{A=1}^{N_n}\dfrac{1}{2M_A}\nabla_{\vec{R}_A}^2-\sum_{i=1}^{N_e}\dfrac{1}{2}\nabla_{\vec{r}_i}^2-\sum_{i=1}^{N_e}\sum_{\mu=1}^{N_n}\dfrac{Z_\mu}{\left|\vec{r}_i-\vec{R}_{\mu}\right|}+\sum_{1\leq i<j\leq N_e}\dfrac{1}{\left|\vec{r}_i-\vec{r}_j\right|}+\sum_{1\leq\nu<\mu\neq N_A}\dfrac{Z_\mu Z_\nu}{\left|\vec{R}_\nu-\vec{R}_\mu\right|}\]
由于原子核的质量相对较大,其动能项较小,可以忽略:
\begin{theorem}[玻恩-奥本海默近似]
	可以将分子的哈密顿量中原子核的动能项忽略,得到
	\[\hat{H}_{\text{BO}}=-\sum_{i=1}^{N_e}\dfrac{1}{2}\nabla_{\vec{r}_i}^2-\sum_{i=1}^{N_e}\sum_{\mu=1}^{N_n}\dfrac{Z_\mu}{\left|\vec{r}_i-\vec{R}_{\mu}\right|}+\sum_{1\leq i<j\leq N_e}\dfrac{1}{\left|\vec{r}_i-\vec{r}_j\right|}+\sum_{1\leq\nu<\mu\neq N_A}\dfrac{Z_\mu Z_\nu}{\left|\vec{R}_\nu-\vec{R}_\mu\right|}\]
\end{theorem}
因此分子的薛定谔方程变为
\[\hat{H}_{\text{BO}}\iPsi=E\iPsi\]
这就是\tbf{BO问题}.
\[\hat{H}_{\text{BO}}(\vec{r};\vec{R})\Phi_k(\vec{r};\vec{R})=\epsilon_k^{\text{BO}}(\vec{R})\Phi_k(\vec{r};\vec{R})\]
其中$\epsilon_k^{\text{BO}}(\vec{R})$即\tbf{绝热势能面(Adiabatic Potential Energy Surfaces)}\footnote{实际上与绝热没有关系,这里的adiabatic更倾向于准静态的意思.}. 势能面对于整个分子的平动,转动具有不变性.\\
\indent 精确地求解分子的薛定谔方程非常困难,因为多个电子和原子会带来相当高的维度,求解也变得非常困难.因此可以先从简单的体系入手寻找规律.
\subsection{氢分子离子}
在BO近似下, \ce{H2+}的薛定谔方程为
\[-\dfrac12\nabla^2\psi+\left(-\dfrac{1}{r_1}-\dfrac{1}{r_2}+\dfrac{1}{R}\right)\psi=E\psi\]
\subsubsection{分离变量法}
处理这一体系可以在长椭球坐标系下采用分离变量法.令
\[\xi=\dfrac{r_1+r_2}{R},\quad\eta=\dfrac{r_1-r_2}{R}\]
以及方位角$\phi$.长椭球坐标系中的Laplace算子为
\[\nabla^2=-\dfrac{4}{R^2(\xi^2-\eta^2)}\left\{\dfrac{\p}{\p\xi}\left[(\xi^2-1)\dfrac{\p}{\p\xi}\right]+\dfrac{\p}{\p\eta}\left[(1-\eta^2)\dfrac{\p}{\p\eta}\right]+\left(\dfrac{1}{\xi^2-1}+\dfrac{1}{1-\eta^2}\dfrac{\p^2}{\p\phi^2}\right)\right\}\]
而
\[\dfrac{1}{r_1}+\dfrac{1}{r_2}=\dfrac{2}{R}\left(\dfrac{1}{\xi+\eta}+\dfrac{1}{\xi-\eta}\right)=\dfrac{4\xi}{R(\xi^2-\eta^2)}\]
令$\tilde{E}=E-\dfrac1R$,代入原方程,在两边乘以$\dfrac14R^2(\xi^2-\eta^2)$可得
\[-\dfrac12\left\{\dfrac{\p}{\p\xi}\left[(\xi^2-1)\dfrac{\p}{\p\xi}\right]+\dfrac{\p}{\p\eta}\left[(1-\eta^2)\dfrac{\p}{\p\eta}\right]+\left(\dfrac{1}{\xi^2-1}+\dfrac{1}{1-\eta^2}\dfrac{\p^2}{\p\phi^2}\right)\right\}\psi-R\xi\psi=\dfrac14R^2(\xi^2-\eta^2)\tilde{E}\psi\]
令$\psi=N(\xi)M(\eta)\dfrac{1}{\sqrt{2\pi}}\e^{\i m\phi}$即可对上面的方程实现变量分离:
\[\left\{\dfrac{\p}{\p\eta}\left[(1-\eta^2)\dfrac{\p}{\p\eta}\right]-\dfrac{m^2}{1-\eta^2}-\dfrac{R^2}{2}\tilde{E}\eta^2-A\right\}M(\eta)=0\]
\[\left\{\dfrac{\p}{\p\xi}\left[(1-\xi^2)\dfrac{\p}{\p\xi}\right]-\dfrac{m^2}{1-\xi^2}-\dfrac{R^2}{2}\tilde{E}\xi^2-2R\xi-A\right\}N(\eta)=0\]
方程中$\tilde{E}$与$A$是待定常数.
\subsubsection{基组展开}
\paragraph{$\symbf\sigma$轨道}
采用\tbf{原子轨道线性组合(Linear Combination of Atomic Orbitals, LACO)},将波函数视作原子轨道的线性组合.这里考虑最小基组,即两个H原子各自的1s轨道,则有
\[\iPsi_{\pm}=C_{\pm}(\iPsi_A\pm\iPsi_B)\]
由此组合的轨道是$\symbf\sigma$\tbf{轨道}.在分离变量的解中, $\sigma$轨道对应的角动量分量$m=0$.从对称性上看, $\sigma$轨道沿键轴方向的对称性与$s$轨道相同,投影图都是圆.
\paragraph{成键与反键轨道}
\indent 现在考虑\tbf{成键轨道}$\iPsi_+$和\tbf{反键轨道}$\iPsi_{-}$.成键轨道的电子密度为
\[\iPsi_+^2\propto\iPsi_A^2+\iPsi_B^2+2\iPsi_A\iPsi_B\]
交叉项导致核间区域干涉相长,电子密度增加.同样地,反键轨道的交叉项导致核间区域干涉相消,电子密度减小.\\
\indent 通过推导可得成键轨道的能量为
\[E_\sigma=E_{\text{H}1\text{s}}+\dfrac{j_0}{R}-\dfrac{j+k}{1+S}\]
其中重叠积分$S$为
\[S=\int\iPsi_A\iPsi_B\d\tau=\left[1+\dfrac{R}{a_0}+\dfrac13\left(\dfrac{R}{a_0}\right)^2\right]\e^{-R/a_0}\]
库伦积分$j$为
\[j=j_0\int\dfrac{\iPsi_A^2}{r_B}\d\tau=\dfrac{j_0}{R}\left[1-\left(1+\dfrac{R}{a_0}\right)\e^{-2R/a_0}\right]\]
交换积分$k$为
\[k=j_0\int\dfrac{\iPsi_A\iPsi_B}{r_B}\d\tau=\dfrac{j_0}{a_0}\left(1+\dfrac{R}{a_0}\right)\e^{-R/a_0}\]
这里$j_0$是能量单位,衡量库伦势能.在原子单位下对于 \ce{H2+}而言$j_0=1$.\\
\indent 反键轨道的电子密度为
\[E_{\sigma^\ast}=E_{\text{H}1\text{s}}+\dfrac{j_0}{R}-\dfrac{j-k}{1-S}\]
反键轨道的能量上升比成键轨道的能量下降的程度更明显.
\subsection{}
\end{document}